%==============================================================================
% Paper D -- skeleton (bullets, to be polished into prose)
% Target venue: signal processing / change-point detection, or a systems
%   track. The team's caution: do NOT claim marker-validated phasic F1 -- the
%   honest contribution is a robust detector PLUS a family-aware validation
%   strategy PLUS the finding that workload markers are not APF-observable.
%
% Build: pdflatex skeleton.tex  (IEEEtran; swap to article if missing).
%==============================================================================
\documentclass[conference]{IEEEtran}
\usepackage{amsmath,booktabs,siunitx,url}

\title{Family-Aware CUSUM Segmentation of Coarse\\
Memory-Liveness Trajectories}

\author{\IEEEauthorblockN{TODO Author list}
\IEEEauthorblockA{TODO Affiliation}}

\begin{document}
\maketitle

\begin{abstract}
% TODO polish. Bullet stub:
\begin{itemize}
  \item Goal: place phase boundaries on a short, coarse APF trajectory and,
        equally, refrain from placing them when the workload is steady.
  \item Method: a windowed CUSUM detector with a median-absolute-deviation
        robust scale and a single drift knob, specialized per workload family.
  \item Finding: steady workloads are correctly left unsegmented (no spurious
        boundaries) for 5 of 6; phasic workloads produce plausible boundary
        counts for 5 of 5 -- but ground-truth marker validation is shown to be
        \emph{unreachable} because the workload markers do not manifest as APF
        mean shifts at this cadence.
\end{itemize}
\end{abstract}

\section{Introduction}
\begin{itemize}
  \item Downstream of capture and tuning, an online monitor wants phase
        boundaries: when does a workload transition between regimes?
  \item The trajectory is hostile to classical detectors: it is short (median 30
        points, range 4--99), coarse (a single scalar per
        $\approx$\SI{10}{\second} median, up to $\approx$\SI{24}{\second}), and
        heavy-tailed.
  \item \textbf{Two failure modes to avoid}: missing real transitions in phasic
        workloads, and hallucinating transitions in steady ones. A good
        segmenter must do both.
  \item \textbf{Approach}: a windowed cumulative-sum (CUSUM) detector with a
        robust (MAD) scale, run on window means; specialized by workload family
        (the family label comes from Paper C's classifier).
  \item \textbf{An honesty-forcing observation}: the synthetic workloads emit
        phase markers, but those markers describe workload-internal events
        (e.g., per-file encryption) that are \emph{not} observable as APF mean
        shifts at this cadence. Marker-aligned F1 is therefore structurally
        unavailable; we validate phasic segmentation by boundary-count
        plausibility instead and report this limitation as a result.
  \item \textbf{Contributions}:
    \begin{itemize}
      \item C1: a windowed CUSUM detector tuned for short, coarse, heavy-tailed
            liveness trajectories (robust scale, single drift knob, explicit
            window-centre-to-snapshot mapping).
      \item C2: a family-aware acceptance scheme -- steady validated by absence
            of spurious boundaries, phasic by plausibility when markers are not
            signal-observable.
      \item C3: an empirical result (10 of 11 workloads accepted) plus the
            documented negative finding that workload markers are not APF
            observables.
    \end{itemize}
\end{itemize}

\section{Detector}
\begin{itemize}
  \item \textbf{Pre-smoothing.} Compute window means
        $\mu[w]$ over the trajectory with window $W$, hop $H$ (length
        $\lfloor (T-W)/H \rfloor + 1$).
  \item \textbf{Robust standardization.} $\sigma = 1.4826 \cdot \mathrm{MAD}(\mu)$
        (fall back to $\mathrm{std}+10^{-9}$ if MAD $=0$);
        $z[t] = (\mu[t] - \mathrm{median}(\mu))/\sigma$.
  \item \textbf{CUSUM recursion.}
        \[
          s^{+}_t = \max(0,\; s^{+}_{t-1} + z[t] - k), \qquad
          s^{-}_t = \min(0,\; s^{-}_{t-1} + z[t] + k).
        \]
        A boundary is emitted when $s^{+}_t > h$ or $s^{-}_t < -h$; both
        accumulators reset to 0 after a detection.
  \item \textbf{Defaults.} $k = 2.0$, $h = 4.0$, minimum separation 2 windows.
  \item \textbf{Mapping.} A detection at window $w$ maps to snapshot index
        $w \cdot H + \lfloor W/2 \rfloor$ (window centre).
  \item \textbf{Family specialization.} Steady workloads keep $k=2.0$; phasic
        workloads use a smaller drift ($k$ down to $\approx 0.1$) to enter the
        detection band on a shallow, slow rhythm.
\end{itemize}

\section{Family-Aware Acceptance}
\begin{itemize}
  \item \textbf{Steady (no boundaries expected).} A cell passes if it is
        stationary and produces few spurious boundaries. Workload-level gate:
        median spurious count $\le 1$ \emph{and} $P(\text{spurious}=0)\ge 0.5$.
  \item \textbf{Phasic, marker-rich (intended path).} Median marker-aligned F1
        $\ge 0.67$ and $P(F1 \ge 0.5) \ge 0.75$. \emph{In this dataset no cell is
        marker-rich} -- see below.
  \item \textbf{Phasic, marker-less (fallback).} Score by plausibility: the
        predicted boundary count must fall in a plausible band
        ($1 \le n_{\text{pred}} \le 20$); workload passes if the plausible
        fraction $\ge 0.70$. \emph{This is a weak gate}: on trajectories of
        $\le 30$ windows a resetting CUSUM lands 1--20 boundaries for nearly any
        non-degenerate input, so the band's acceptance probability under a
        stationary null is close to 1. ``5 of 5 phasic pass'' therefore carries
        little evidential weight; the band must be calibrated against a
        permutation / stationary-null distribution and its false-accept rate
        reported (see Limitations).
  \item \textbf{Regression guard (G4).} Placeholder in this version (always
        passes); wiring it to a legacy baseline is future work.
\end{itemize}

\section{Results}
\begin{itemize}
  \item \textbf{Overall.} 10 of 11 workloads pass acceptance
        (120 of 132 cells). Detector configuration: $k=2.0$, $h=4.0$,
        plausibility mode for phasic.
  \item \textbf{Steady (5 of 6 pass).} Stationary with zero spurious boundaries:
        \texttt{hashtable} ($P_0=0.83$), \texttt{rmw} (0.83),
        \texttt{workingset} (0.75), \texttt{mmap} (0.67), \texttt{writemag}
        (0.50). \textbf{Failure}: \texttt{mem\_pagefault\_density\_v2} --
        median stationarity 0.67, median spurious 1.0, $P_0=0.33$ -- it injects
        spurious boundaries and fails the steady gate.
  \item \textbf{Phasic (5 of 5 pass, by plausibility).} All 12 cells per phasic
        workload are marker-less; plausible-fraction:
        \texttt{scanner} 1.00, \texttt{slowburn} 0.92, \texttt{batched} 0.83,
        \texttt{seq} 0.75, \texttt{selective} 0.75 -- each $\ge 0.70$.
  \item \textbf{The negative finding (marker F1).} The marker-rich gate G1 has
        \emph{zero} applicable cells: the workload markers do not appear as APF
        mean shifts at this cadence (compounded by a capture-side wall-clock
        skew), so marker-aligned F1 is unreachable. This is a statement about
        \emph{observability}, not detector quality, and it motivates either
        re-instrumentation or a finer signal.
  \item TODO: figure -- CUSUM accumulators over an example phasic trajectory;
        table -- per-workload (family, $W/H$, stationarity, spurious, plausible
        fraction, verdict).
\end{itemize}

\section{Limitations}
\begin{itemize}
  \item Phasic acceptance rests on plausibility, not ground-truth F1; the
        strong claim (``we segment phases correctly'') is not yet supported and
        is explicitly deferred. The plausibility band ($1 \le n \le 20$) is
        moreover near-unfalsifiable at these trajectory lengths and needs a
        null-calibrated false-accept rate before it counts as validation.
  \item Short trajectories give CUSUM little warmup; the robust scale is
        estimated from few windows.
  \item The steady failure (\texttt{pagefault}) and the placeholder regression
        guard (G4) are open items.
\end{itemize}

\section{Conclusion}
\begin{itemize}
  \item A robust, family-aware CUSUM cleanly distinguishes ``leave it alone''
        (steady) from ``segment it'' (phasic) for 10 of 11 workloads.
  \item The honest boundary of the result is observability: validating phasic
        boundaries against ground truth needs markers that survive into the APF
        signal -- the central lesson for the next capture iteration.
\end{itemize}

\bibliographystyle{IEEEtran}
\bibliography{../refs}
\end{document}
